\documentclass[11pt,a4paper]{article}

\usepackage{fontspec}
\setmainfont{DejaVu Sans}
\setsansfont{DejaVu Sans}
\setmonofont{DejaVu Sans Mono}[Scale=0.88]
\defaultfontfeatures{Ligatures=TeX}
\usepackage{microtype}
\usepackage{unicode-math}
\setmathfont{Latin Modern Math}[Scale=1.0]
\usepackage[ngerman]{babel}
\usepackage{geometry}
\geometry{margin=2.5cm}

% --- Zeilenumbruch: sauber wie im Harmonikabuch ---
\sloppy
\emergencystretch=3em
\tolerance=2500
\hbadness=10000
\hfuzz=2pt
\clubpenalty=10000
\widowpenalty=10000

\usepackage{amsmath,amssymb}
\usepackage{parskip}
\usepackage{booktabs}
\usepackage{array}
\usepackage{xcolor}
\usepackage{hyperref}
\hypersetup{
  colorlinks=true,
  linkcolor=blue!50!black,
  urlcolor=blue!50!black,
  pdftitle={Kosmologische Rotverschiebung in FFGFT - Kurzfassung},
  pdfauthor={Johann Pascher}
}

\title{Kosmologische Rotverschiebung in FFGFT\\[2pt]
  \large Kurzfassung für die IPI-Korrespondenz}
\author{Johann Pascher\\[2pt]
\small ORCID: \href{https://orcid.org/0009-0000-6518-4064}{0009-0000-6518-4064}\\
\small GitHub: \url{https://github.com/jpascher/T0-Time-Mass-Duality}\\
\small Zenodo (FFGFT v1.0.13): \url{https://zenodo.org/records/20041529}\\
\small DOI: \href{https://doi.org/10.5281/zenodo.20041529}{10.5281/zenodo.20041529}}
\date{Mai 2026}

\begin{document}
\maketitle

\begin{abstract}
\noindent In FFGFT ist die kosmologische Rotverschiebung ein geometrischer Pfad-Effekt, keine Doppler-Signatur einer Raumzeit-Expansion. Photonen durchqueren das fluktuierende $\xi$-Vakuum entlang Geodäten, deren effektive Pfadlänge $L_{\mathrm{eff}}$ systematisch länger ist als die Euklidische Distanz $d$, mit $z = (L_{\mathrm{eff}} - d)/d$. Die Hubble-Konstante wird zu einer fundamentalen Kopplung, die aus dem einzigen dimensionslosen $\xi = 4/30{,}000$ folgt:
\[
\frac{\hbar H_0}{m_e c^2} = \frac{\pi}{2}\,\xi^{10} \qquad\Longrightarrow\qquad H_0 \approx 66{,}5\ \text{km/s/Mpc},
\]
parameterfrei, innerhalb $1{,}3\%$ des Planck-Werts. Diese Notiz fasst den FFGFT-Rotverschiebungsmechanismus, die zwei äquivalenten Herleitungen von $H_0$ und die natürliche Auflösung der Hubble-Spannung zusammen. Quelldokumente im FFGFT-Korpus: Doc~025 (Kosmologie), Doc~026 (geometrische Kosmologie), Doc~027 (Sonnensystem-Konsistenz), Doc~165 ($H_0$-Verhältnis), Doc~166 (Casimir-Hubble-Brücke).
\end{abstract}

\section{Ausgangspunkt}

FFGFT modelliert das Universum als statisch und flach. Der einzige dimensionslose Parameter $\xi = 4/30{,}000$ regelt die Granularität des $\xi$-Vakuums. Photonen, die das Vakuum durchqueren, werden nicht von einer expandierenden Metrik getragen, sondern propagieren entlang Geodäten, die die kleinskaligen Fluktuationen des $\xi$-Felds inkorporieren. Der kumulative Effekt über kosmologische Distanzen ist eine systematische Rotverschiebung; über Sonnensystem-Distanzen ist er operational vernachlässigbar.

\section{Rotverschiebung als geometrischer Pfad-Effekt}

Die zentrale Beobachtung, formalisiert in Doc~026 über eine Finite-Elemente-Simulation des $\xi$-Vakuums, ist, dass Geodäten durch das körnige $\xi$-Feld keine perfekten Geraden sind. Die effektive Pfadlänge $L_{\mathrm{eff}}$ ist systematisch größer als die direkte Euklidische Distanz $d$ zwischen Quelle und Beobachter:
\begin{equation}
z = \frac{L_{\mathrm{eff}} - d}{d}\,,
\label{eq:z_geometric}
\end{equation}
mit der linearisierten Approximation
\begin{equation}
z \approx d \cdot C \cdot \xi
\label{eq:z_linear}
\end{equation}
wobei $C$ eine geometrische Konstante ist und $\xi$ die kumulative Pfadstreckungsrate steuert. Die Rotverschiebung ist eine Eigenschaft der Raumzeit-Geometrie selbst, nicht der Relativgeschwindigkeit zwischen Quelle und Beobachter.

\paragraph{Frequenzunabhängigkeit als struktureller Test.}
Eine Geodäte ist eine Eigenschaft der Geometrie; sie hängt nicht von der Photonenfrequenz ab. Ein rotes und ein blaues Photon, die von derselben Quelle ausgesandt werden, folgen \emph{demselben} verlängerten Pfad. Ihre Wellenlängen werden um denselben prozentualen Faktor gestreckt. Dies stimmt mit der beobachteten Frequenzunabhängigkeit der kosmologischen Rotverschiebung überein -- ein Merkmal, das konventionelle ``Tired Light''-Modelle nicht reproduzieren können, und das den geometrischen FFGFT-Mechanismus von Energieverlust-Szenarien mit frequenzabhängigem Wirkungsquerschnitt unterscheidet.

\section{Die Hubble-Konstante als fundamentales \texorpdfstring{$\xi$}{xi}-Verhältnis}

Im Korpus sind zwei äquivalente FFGFT-Herleitungen von $H_0$ aufgezeichnet.

\paragraph{Doc~165 -- das dimensionslose Verhältnis.} Die sauberste Formulierung drückt $H_0$ als reines Verhältnis zwischen kosmologischer und Elektronen-Ruheenergieskala aus:
\begin{equation}
\boxed{\;\frac{\hbar H_0}{m_e c^2} = \frac{\pi}{2}\cdot\xi^{10}\;}
\label{eq:h0_ratio}
\end{equation}
Numerisch:
\begin{equation}
H_0^{\mathrm{FFGFT}} \approx 66{,}5\ \text{km/s/Mpc} \qquad (\text{Planck: } 67{,}4 \pm 0{,}5,\ \text{Abweichung } +1{,}3\%).
\end{equation}
Der ganzzahlige Exponent $10$ spiegelt wider, wie weit die kosmologische Energieskala von der Lepton-Massenskala entfernt liegt, gemessen in $\xi$-Einheiten. Der geometrische Vorfaktor $\pi/2$ entstammt der Torus-Topologie des $\xi$-Zeitfelds.

\paragraph{Doc~026 -- die Hubble-Energie-Form.} Äquivalent wird die Hubble-Energie ausgedrückt als
\begin{equation}
E_H = E_0\cdot\xi^{41/4} = 1{,}41\times 10^{-33}\,\text{eV}\,, \qquad H_0 = \frac{E_H}{\hbar} \approx 66{,}2\ \text{km/s/Mpc}\,,
\label{eq:EH_026}
\end{equation}
mit $E_0 = 7{,}398$~MeV als FFGFT-Referenzenergie. Der Exponent zerlegt sich als
\begin{equation}
\frac{41}{4} = \underbrace{10}_{\text{Doc 165 (}H_0\text{-Verhältnis)}} + \underbrace{\frac{1}{4}}_{\text{Doc 091 (Casimir-Skala)}}\,,
\end{equation}
wodurch die Hubble-Herleitung mit der Casimir-CMB-Herleitung über ein und dasselbe $\xi$ verknüpft wird.

\paragraph{Reinterpretation.} Beide Formen stimmen miteinander und mit der Planck-Skalen-Messung auf besser als $1\%$ überein. In FFGFT ist $H_0$ keine Expansionsrate. Es ist eine fundamentale Kopplung, die die Stärke der Photon-Vakuum-Wechselwirkung beschreibt, abgeleitet allein aus $\xi$.

\section{Die Hubble-Spannung}

Die Standardkosmologie sieht sich einer $\sim 9\%$-Diskrepanz zwischen Frühuniversum-Messungen ($H_0 \approx 67{,}4$ km/s/Mpc, Planck CMB) und Spätuniversum-Messungen ($H_0 \approx 73$ km/s/Mpc, SH0ES-Distanzleiter; H0LiCOW-Linsen) gegenüber. Innerhalb von $\Lambda$CDM ist dies ein strukturelles Problem, weil die Expansionsrate eine einzige universelle Konstante sein sollte.

In FFGFT löst sich die Spannung auf konzeptioneller Ebene aus drei Gründen auf.

\begin{enumerate}\itemsep2pt
\item \textbf{Keine Expansion, keine Expansionsrate.} $H_0$ ist nicht die Rate, mit der sich der Raum ausdehnt. Sie quantifiziert die kumulative Photon-Vakuum-Wechselwirkung. Verschiedene Messmethoden müssen nicht identische Werte liefern, und die Diskrepanz zwischen Frühuniversum- und Spätuniversum-Inferenzen ist für sich genommen kein Widerspruch.

\item \textbf{Methodenabhängige Pfad-Sondierungen.} Verschiedene Sondierungen (CMB-akustische Peaks, Distanzleiter-Cepheiden, Gravitationslinsen-Zeitverzögerungen) sondieren verschiedene Distanzen und verschiedene Photonen-Energiebereiche. In einem körnigen $\xi$-Vakuum kann die in jedem Regime effektive kumulative Kopplung systematisch unterschiedlich sein. Das Muster früh-niedrig / spät-hoch ist genau das, was man erwarten würde, wenn die $\xi$-Kopplung mild von den integrierten Pfadeigenschaften abhängt.

\item \textbf{Anisotropie-Verträglichkeit.} Finite-Elemente-Simulationen des $\xi$-Vakuums (Doc~026) liefern kleine richtungsabhängige Variationen der effektiven Kopplung, in Übereinstimmung mit berichteten Hinweisen auf großräumige Anisotropie in $H_0$-Messungen.
\end{enumerate}

\section{Konsistenz mit Sonnensystem-Daten}

Eine natürliche Sorge bei jeder Modifikation der kosmologischen Rotverschiebung ist die Verträglichkeit mit hochpräzisen Sonnensystem-Tests. FFGFT adressiert dies direkt (Doc~027): Die geometrische Pfadstreckung ist ein kumulativer Effekt über kosmologische Distanzen. Über Sonnensystem-Distanzen (Lichtstunden) ist die Streckung operational vernachlässigbar, und FFGFT sagt \emph{null} messbare Bahnanomalien im Sonnensystem voraus.

Dies ist ein struktureller Unterschied zu modifizierten-Gravitations-Ansätzen für die Hubble-Spannung, von denen mehrere von Nathan et al.~(2024, MNRAS 544, 975) auf Sonnensystem-Grundlage ausgeschlossen wurden. FFGFT modifiziert die Gravitation auf Sonnensystem-Skalen nicht; es reinterpretiert die Prämisse der kosmologischen Rotverschiebung. Die Cassini-Mission-Beschränkungen, die viele modifizierte-Gravitations-Vorschläge ausschließen, sind daher automatisch erfüllt.

\section{Zusammenfassung der Vorhersagen}

\begin{enumerate}\itemsep2pt
\item Kosmologische Rotverschiebung ist ein geometrischer Pfad-Effekt, $z = (L_{\mathrm{eff}}-d)/d$, frequenzunabhängig per Konstruktion.
\item $H_0$ folgt parameterfrei aus $\xi = 4/30{,}000$:
\[
\hbar H_0 / (m_e c^2) = (\pi/2)\,\xi^{10} \quad\Longleftrightarrow\quad E_H = E_0\,\xi^{41/4}\,,
\]
mit $H_0 \approx 66{,}2$ -- $66{,}5$ km/s/Mpc.
\item Das Hubble-Spannungs-Muster (früh-niedrig / spät-hoch) ist mit FFGFT strukturell verträglich und erfordert keine zusätzlichen Modellkomponenten.
\item Die Cassini-Sonnensystem-Beschränkungen sind automatisch erfüllt; keine Modifikation der Gravitation auf lokalen Skalen.
\item Die Casimir-Skala $L_\xi$ und die Hubble-Konstante $H_0$ erfüllen eine algebraische Brückenrelation, $L_\xi \cdot H_0/c \propto \xi^{41/4}$ (Doc~166).
\end{enumerate}

\section*{Dokumentenverweise}

\noindent\small\begin{tabular}{ll}
\textbf{Doc 025} & \emph{T0-Kosmologie.} Statisches, ewiges Universum; CMB als $\xi$-Feld-Manifestation; \\
& drei alternative Interpretationen des Rotverschiebungsmechanismus \\
& (Energieverlust, kumulative Gravitation, Raumzeit-Geometrie). \\[2pt]
\textbf{Doc 026} & \emph{Geometrische Kosmologie.} Finite-Elemente-Simulation des $\xi$-Vakuums; \\
& Rotverschiebung als Pfadstreckung; Herleitung von $E_H = E_0\,\xi^{41/4}$; \\
& dreistufige Auflösung der Hubble-Spannung. \\[2pt]
\textbf{Doc 027} & \emph{MNRAS-Analyse.} Verträglichkeit mit Sonnensystem-Tests; \\
& Vergleich mit Nathan et al.\ (2024) zur modifizierten Gravitation. \\[2pt]
\textbf{Doc 165} & \emph{Das Hubble-Verhältnis.} Herleitung von $\hbar H_0/(m_e c^2) = (\pi/2)\,\xi^{10}$ \\
& mit $H_0 \approx 66{,}5$ km/s/Mpc, parameterfrei. \\[2pt]
\textbf{Doc 166} & \emph{Casimir-Hubble-Brücke.} Die Zerlegung $41/4 = 10 + 1/4$ \\
& und die algebraische Beziehung $L_\xi \cdot H_0/c$. \\
\end{tabular}

\medskip
Alle Dokumente verfügbar unter \url{https://github.com/jpascher/T0-Time-Mass-Duality}\\
und als Zenodo-Release v1.0.13: \url{https://zenodo.org/records/20041529}.

\end{document}
